\section*{CRediT authorship contribution statement}

Jingquan Wang: Conceptualization, Methodology, Software, Validation, Formal analysis, Writing --
original draft, Writing -- review and editing.

\section*{Declaration of competing interest}

The author declares no competing interests relevant to this work.

\section*{Funding}

No external funding was received for this research.

\section*{Data availability}

The implementation, the experiment scripts that generate every table and figure of
Section~\ref{sec:numerical}, the result files, and the Lean~4 development are available in the
repository described in \ref{app:repro}.

\section*{Declaration of generative AI and AI-assisted technologies in the writing process}

During preparation of this work, the author used an AI coding assistant for drafting support,
consistency checking and script editing. The author reviewed and edited the content and takes full
responsibility for the published article.

\appendix

\section{Lower-pair constraint rows of the benchmark mechanisms}
\label{app:lower-pair}

The four benchmark mechanisms use the scalar constraint rows of the public suite
\cite{kissel2022absolute}. Let body $i$ have centre position $r_i$, rotation matrix $A_i\in SO(3)$,
translational velocity $v_i$, angular velocity $\omega_i$, translational acceleration $a_i$ and
angular acceleration $\alpha_i$. A marked point and a marked direction fixed in body $i$ are
$x_i=r_i+A_i s_i$ and $u_i=A_i a_i$. The row families are
\begin{align}
  C_{\mathrm{CD},k}(q,t) &= e_k^T\left(x_j-x_i\right)-d_k(t), &
  C_{\mathrm{DP1}}(q,t) &= u_i^T u_j-c(t),\\
  C_{\mathrm{DP2}}(q,t) &= u_i^T\left(x_j-x_i\right)-d(t), &
  C_{\mathrm{D}}(q,t) &= \left\|x_j-x_i\right\|^2-\ell^2(t),
  \label{eq:lower-pair-row-families}
\end{align}
where CD is a coordinate difference, DP1 constrains two body-fixed directions, DP2 constrains a
body-fixed direction against a point difference, and D is a distance row. Revolute joints are
assembled from CD and DP1 rows; the slider-crank prismatic joint additionally uses DP2 and distance
rows; drivers prescribe $c(t)$ of a DP1 row. With
$\dot x_i = v_i+\omega_i\times A_i s_i$, $\dot u_i = \omega_i\times u_i$,
$\ddot x_i = a_i+\alpha_i\times A_i s_i+\omega_i\times(\omega_i\times A_i s_i)$,
$\ddot u_i = \alpha_i\times u_i+\omega_i\times(\omega_i\times u_i)$, $\Delta x=x_j-x_i$ and
its derivatives, the velocity- and acceleration-level rows enforced at every stage are
\begin{align}
  \dot C_{\mathrm{CD},k} &= e_k^T\Delta\dot x-\dot d_k, &
  \ddot C_{\mathrm{CD},k} &= e_k^T\Delta\ddot x-\ddot d_k,\\
  \dot C_{\mathrm{DP1}} &= \dot u_i^T u_j+u_i^T\dot u_j-\dot c, &
  \ddot C_{\mathrm{DP1}} &= \ddot u_i^T u_j+2\dot u_i^T\dot u_j+u_i^T\ddot u_j-\ddot c,\\
  \dot C_{\mathrm{DP2}} &= \dot u_i^T\Delta x+u_i^T\Delta\dot x-\dot d, &
  \ddot C_{\mathrm{DP2}} &= \ddot u_i^T\Delta x+2\dot u_i^T\Delta\dot x+u_i^T\Delta\ddot x-\ddot d,\\
  \dot C_{\mathrm{D}} &= 2\Delta x^T\Delta\dot x-2\ell\dot\ell, &
  \ddot C_{\mathrm{D}} &= 2\Delta\dot x^T\Delta\dot x+2\Delta x^T\Delta\ddot x-2(\dot\ell^2+\ell\ddot\ell).
  \label{eq:lower-pair-acceleration-rows}
\end{align}

\paragraph{Mechanism data}
All bodies are steel bars of square cross-section $0.05$~m and density $7800$~kg/m$^3$ unless
stated otherwise; $e_x,e_y,e_z$ are the coordinate directions and body-fixed vectors are given in
the body frame.

\emph{Single pendulum.} One bar of length $4$~m (mass $78$~kg,
$J=\operatorname{diag}(0.0325,104.016,104.016)$~kg\,m$^2$). Ground revolute joint: CD rows with
$s_1=(-2,0,0)$, $s_0=0$; DP1 rows $(e_x,e_x)$, $(e_y,e_x)$. Driven DP1 row $(e_y,-e_z)$ with
$c(t)=\cos\theta(t)$, $\theta(t)=\pi/2+(\pi/4)\cos 2t$.

\emph{Double pendulum.} Bars of length $4$ and $2$~m (masses $78$ and $39$~kg). Ground revolute
joint as above; inter-link revolute joint: CD rows with $s_1=(2,0,0)$, $s_2=(-1,0,0)$, DP1 rows on
body 2 against ground $(e_x,e_x)$, $(e_y,e_x)$. Released from the horizontal configuration
$(\theta_1,\theta_2)(0)=(-\pi/2,0)$ at rest; no friction, no drive.

\emph{Four-link loop.} Rotor, link 2 and link 3 with masses $2,1,1$~kg and inertia diagonals
$(4,2,0)$, $(12.4,0.01,0)$, $(4.54,0.01,0)$~kg\,m$^2$. Revolute joint A (body 1--ground): CD
$s_1=s_0=0$, DP1 $(e_x,e_x)$, $(e_y,e_x)$. Revolute joint D (body 3--ground): CD $s_3=(0,3.7,0)$,
$s_0=(-4,-8.5,0)$, DP1 $((0,0,1),e_y)$, $(e_y,e_y)$. Spherical joint (bodies 3--2): CD
$s_3=(0,-3.7,0)$, $s_2=(0,-6.1,0)$. Universal joint (bodies 1--2): CD $s_1=(-2,0,0)$,
$s_2=(0,6.1,0)$, DP1 $((0,-0.662,0.75),e_z)$. Driving DP1 row $((-1,0,0),e_y)$ with
$c(t)=\cos(\pi t+\pi/2)$.

\emph{Slider-crank.} Crank, rod and slider with masses $0.12,0.5,2$~kg and inertia diagonals
$(10^{-4},10^{-5},10^{-4})$, $(4\times10^{-3},4\times10^{-4},4\times10^{-3})$,
$(10^{-4},10^{-4},10^{-4})$~kg\,m$^2$. Revolute joint A (ground--crank): CD $s_4=(0,0.1,0.12)$,
$s_1=0$, DP1 $(e_y,e_x)$, $(e_z,e_x)$. Spherical joint B (rod--crank): CD $s_2=(-0.15,0,0)$,
$s_1=(0,0.08,0)$. Joint C (rod--slider): DP1 $(e_y,e_z)$ and two DP2 rows from rod point
$(0.15,0,0)$ to slider point $0$ along $e_y$, $e_z$. Translational joint D (slider--ground): DP1
$(-e_x,e_x)$, $(e_y,e_x)$, $(-e_x,e_y)$ and DP2 rows from slider point $0$ to ground point
$(0.2,0,0)$ along $-e_x$, $e_y$. Distance row between rod point $(0.15,0,0)$ and slider point
$(0,0,1)$ with $\ell=1$. Driving DP1 row $(e_y,e_y)$ with $c(t)=\cos(-2\pi t+\pi/2)$.

\section{Diagnostics of the analysis constants}
\label{app:diagnostics}

Three diagnostics of Section~\ref{sec:order} are evaluated on the chain benchmark ($v_s=0.5$) over the
short window $[0,0.08]$ with $h=0.04$, $0.02$, $0.01$ and a tight Newton tolerance of $10^{-13}$, so that
the quantities measured are those of the converged stage vectors and not of the solver; the solver
envelope (Table~\ref{tab:e7}) is measured on the sweep of Section~\ref{sec:e1} itself.
(i)~Each converged stage vector is read in the
joint-coordinate chart and the reduced collocation defects $\Delta_i[s_j]$, $\Delta_i[\dot s_j]$,
$\Delta_i[\theta_j]$, $\Delta_i[\dot\theta_j]$ of Eq.~\eqref{eq:collocation-defect-operator} are
evaluated together with the component of $d_j,\dot d_j,\ddot d_j$ orthogonal to $n_j$: all are at
roundoff (Table~\ref{tab:exact-identity-check}), as Lemma~\ref{lem:exact-stage-identity} states.
(ii)~The constants of Lemma~\ref{lem:p2-from-p1} are measured from the implemented Jacobian at the
converged stage and at the root of the $h=0$ stage system (Table~\ref{tab:p2-constants-check}): the
conclusion $\|J_h^{-1}\|\le2\|J_0^{-1}\|$ holds with ratio at most $1.55$, while the sufficient
Neumann condition would need $h\le2\times10^{-5}$ in the Euclidean norm and $h\le8\times10^{-4}$ in
the endpoint-linearized norm; without joint friction the threshold rises to $1.4\times10^{-2}$.
(iii)~The distances of the two predictors to $Z_G$ (Table~\ref{tab:predictor-check}) confirm the
discussion after Lemma~\ref{lem:newton-envelope}: the lifted endpoint predictor is $O(h)$ from the
Gauss stage, the implemented predictor is at distance about $50$ and its first Newton step expands
that distance by up to $18\%$.

\begin{table}[!htb]
\centering
\caption{Exact stage identity at the converged stages of the chain (Newton tolerance $10^{-13}$,
window $[0,0.08]$): maximal implemented residual, maximal reduced collocation defect and maximal
component of $d_j,\dot d_j,\ddot d_j$ orthogonal to $n_j$.}
\label{tab:exact-identity-check}
\small
\begin{tabular}{P{0.07\linewidth}P{0.08\linewidth}P{0.13\linewidth}P{0.17\linewidth}P{0.17\linewidth}P{0.17\linewidth}}
\toprule
$h$ & steps & Newton total & max stage residual & max reduced defect & max $\perp n_j$ component \\
\midrule
0.04 & 2 & 12 & $6.8\times10^{-14}$ & $9.0\times10^{-15}$ & $1.5\times10^{-15}$ \\
0.02 & 4 & 24 & $1.5\times10^{-14}$ & $5.5\times10^{-16}$ & $1.9\times10^{-15}$ \\
0.01 & 8 & 48 & $1.6\times10^{-14}$ & $5.8\times10^{-16}$ & $2.0\times10^{-15}$ \\
\bottomrule
\end{tabular}
\end{table}

\begin{table}[!htb]
\centering
\caption{Constants of Lemma~\ref{lem:p2-from-p1} on the chain (window $[0,0.08]$): $M_0=\max_n\|J_0^{-1}\|$,
$C_J=\max_n\|J_h-J_0\|/h$, observed inverse ratio $\|J_h^{-1}\|/\|J_0^{-1}\|$ and Neumann threshold
$h_0=1/(2M_0C_J)$, in the Euclidean norm on the stage layout (E) and in the endpoint-linearized
norm $\|J_0^{-1}\,\cdot\,\|$ (L, where $M_0=1$); the last two columns repeat the L-norm quantities
for the same mechanism without joint friction.}
\label{tab:p2-constants-check}
\scriptsize
\setlength{\tabcolsep}{3pt}
\begin{tabular}{@{}lrlrllllr@{}}
\toprule
$h$ & $M_0$ (E) & $C_J$ (E) & ratio & $h_0$ (E) & $C_J$ (L) & $h_0$ (L) & $C_J$ (L, no fr.) & $h_0$ (L, no fr.) \\
\midrule
0.04 & 18.2 & $7.5\mathrm{e}{2}$ & 1.55 & $3.7\mathrm{e}{-5}$ & $5.1\mathrm{e}{2}$ & $9.8\mathrm{e}{-4}$ & 34.1 & $1.5\mathrm{e}{-2}$ \\
0.02 & 24.3 & $8.2\mathrm{e}{2}$ & 1.26 & $2.5\mathrm{e}{-5}$ & $5.6\mathrm{e}{2}$ & $8.9\mathrm{e}{-4}$ & 35.5 & $1.4\mathrm{e}{-2}$ \\
0.01 & 28.4 & $9.6\mathrm{e}{2}$ & 1.18 & $1.8\mathrm{e}{-5}$ & $6.5\mathrm{e}{2}$ & $7.6\mathrm{e}{-4}$ & 36.5 & $1.4\mathrm{e}{-2}$ \\
\bottomrule
\end{tabular}
\end{table}

\begin{table}[!htb]
\centering
\caption{Predictor distances for Lemma~\ref{lem:newton-envelope} on the chain (window $[0,0.08]$,
Euclidean norm on the stage layout): distance of the lifted endpoint predictor $Z^{(0)}_\ast$ and of the implemented
predictor to $Z_G$, and the ratio $\|Z^{(1)}-Z_G\|/\|Z^{(0)}-Z_G\|$ of the first full Newton step
from the implemented predictor.}
\label{tab:predictor-check}
\small
\begin{tabular}{P{0.08\linewidth}P{0.20\linewidth}P{0.22\linewidth}P{0.20\linewidth}}
\toprule
$h$ & lifted predictor & implemented predictor & first-step ratio \\
\midrule
0.04 & 3.23 & 50.2 & 1.12 \\
0.02 & 2.11 & 50.3 & 1.16 \\
0.01 & 1.07 & 50.4 & 1.18 \\
\bottomrule
\end{tabular}
\end{table}

\paragraph{Solver envelope}
Table~\ref{tab:e7} records, for every grid of the sweep, the largest residual of the accepted Newton
iterate under the production stopping rule and its ratio to $h^7$, which is the constant $c_\eta$
with which that grid satisfies (H3). The accepted residual is set by the tolerance, not by $h$, so
the ratio grows from $10^{-4}$ to $10^{6}$ across the sweep; as discussed after
Lemma~\ref{lem:newton-envelope}, the resulting bound $2M\|F_h(\widetilde Z)\|$ on the stage
perturbation is pessimistic by three orders of magnitude relative to the observed errors. The
endpoint position-level constraint norm of Table~\ref{tab:e1} is the quantity that actually
dominates the deviation of the numerical trajectory from the constraint manifold; it decreases as
$h^6$ as the theorem requires.

\begin{table}[!htb]
\centering
\caption{Solver envelope on the chain sweep ($\tau_N=10^{-11}$): largest residual norm of the
accepted iterate over all steps, its ratio to $h^7$ (the $c_\eta$ of (H3) for that grid), and the
largest components of the stage displacements $d_j,\dot d_j,\ddot d_j$ orthogonal to $n_j$.}
\label{tab:e7}
\footnotesize
\begin{tabular}{@{}lrlll@{}}
\toprule
$h$ & Newton/step & $\max_n\|F_h(\widetilde Z_n)\|$ & ratio to $h^7$ ($c_\eta$) & max $\perp n_j$ component \\
\midrule
$0.1$      & 6.00 & $8.1\mathrm{e}{-12}$ & $8.1\mathrm{e}{-5}$ & $5.6\mathrm{e}{-13}$ \\
$0.05$     & 6.00 & $3.5\mathrm{e}{-13}$ & $4.5\mathrm{e}{-4}$ & $3.2\mathrm{e}{-15}$ \\
$0.025$    & 5.90 & $8.8\mathrm{e}{-12}$ & $1.4$               & $6.4\mathrm{e}{-14}$ \\
$0.0125$   & 5.72 & $2.7\mathrm{e}{-12}$ & $5.7\mathrm{e}{1}$  & $4.3\mathrm{e}{-15}$ \\
$0.00625$  & 5.72 & $3.1\mathrm{e}{-12}$ & $8.3\mathrm{e}{3}$  & $4.1\mathrm{e}{-15}$ \\
$0.003125$ & 5.72 & $3.4\mathrm{e}{-12}$ & $1.2\mathrm{e}{6}$  & $5.3\mathrm{e}{-15}$ \\
\bottomrule
\end{tabular}
\end{table}

\section{Lean development}
\label{app:lean}

The Lean~4 development (Mathlib pinned by its manifest) machine-checks the algebraic and
perturbation lemmas of Section~\ref{sec:order}: $57$ theorems, no open goals, and an exhaustive
axiom audit showing that every theorem depends only on the standard axioms
\nolinkurl{propext}, \nolinkurl{Classical.choice} and \nolinkurl{Quot.sound}. It builds with
\texttt{lake build}; the script \nolinkurl{scripts/check.sh} runs the build, the axiom audit and
a lint pass. Table~\ref{tab:lean-development} maps the statements of the paper to the theorems that
check them. The row identities are proved for a transcription of the implemented residual into
Lean; the perturbation lemmas are proved in the abstract normed-space setting with the constants of
the paper. Butcher's order theorem and the Gauss local-error estimate are not formalized and enter
as cited results.

\begin{table}[!htb]
\centering
\caption{Statements of the paper and the Lean theorems that check them.}
\label{tab:lean-development}
\footnotesize
\setlength{\tabcolsep}{4pt}
\begin{tabular}{P{0.40\linewidth}P{0.52\linewidth}}
\toprule
Statement & Lean theorem \\
\midrule
Lemma~\ref{lem:exact-stage-identity}, non-dynamic rows (both directions) &
\nolinkurl{FullVA.Transition.nondynamic_rows_iff} \\
Lemma~\ref{lem:exact-stage-identity}, Newton--Euler rows &
\nolinkurl{NewtonEuler.dynamic_rows_vanish} \\
Lemma~\ref{lem:p2-from-p1}, bound Eq.~\eqref{eq:p2-perturbation-bound} &
\nolinkurl{uniform_inverse_of_perturbation} \\
Lemma~\ref{lem:stage-residual-defect}, root uniqueness &
\nolinkurl{root_unique_of_linearization} \\
Lemma~\ref{lem:gauss-endpoint-defect}, bound Eq.~\eqref{eq:gauss-quadrature-defect} &
\nolinkurl{gauss6_step_defect} \\
Butcher conditions $B(6)$, $C(3)$, $D(3)$; failure of $B(7)$ &
\nolinkurl{Gauss6.butcher_order_six_hypotheses}, \nolinkurl{Gauss6.not_B_seven} \\
Lemma~\ref{lem:endpoint-closure} &
\nolinkurl{endpoint_closure_exists}, \nolinkurl{endpoint_correction_bound_h7} \\
Lemma~\ref{lem:inexact-newton} &
\nolinkurl{inexact_newton_output_bound} \\
Lemma~\ref{lem:newton-envelope}, decay Eq.~\eqref{eq:newton-envelope-decay} &
\nolinkurl{simplified_newton_residual_decay} \\
Lemma~\ref{lem:local-global-transfer} &
\nolinkurl{local_to_global}, \nolinkurl{reported_grid_bound} \\
Theorem~\ref{thm:g6fullva-order}, three-term chain and grid bounds &
\nolinkurl{local_defect_bound}, \nolinkurl{conditional_sixth_order_grid_bound} \\
\bottomrule
\end{tabular}
\end{table}

\section{Reproducibility}
\label{app:repro}

The repository accompanying the paper has five top-level folders: \texttt{paper/} (this
manuscript), \texttt{proof/} (the Lean development), \texttt{numerics/} (the implementation and all
experiment scripts), \texttt{external/} (the public benchmark codes) and \texttt{validation/} (the
audit trail of the development).
\begin{itemize}
  \item The step for the chain is implemented by the single residual function in
  \texttt{numerics/v047\_cylindrical\_chain\_pipeline/run\_v047.py}; its automatic derivative is the
  Newton Jacobian. The assemblies for the benchmark mechanisms are referenced from the experiment
  scripts below.
  \item The experiments of Section~\ref{sec:numerical} are the scripts \texttt{e1\_convergence.py}
  to \texttt{e7\_solver\_envelope.py} in \texttt{numerics/v049\_paper\_experiments/}; each writes a
  CSV and a JSON file under \texttt{results/}, and \texttt{make\_figures.py} draws the figures from
  those files.
  \item \texttt{validation/validate\_manuscript.py} checks that every number in the tables of
  Section~\ref{sec:numerical} and of \ref{app:diagnostics} equals the value in its result file, that
  the LaTeX log is clean, and that the theorem names of Table~\ref{tab:lean-development} exist in
  \texttt{proof/}.
  \item The Lean development is checked by \texttt{proof/scripts/check.sh} (build, exhaustive axiom
  audit, linter).
\end{itemize}
All runs reported were made on one CPU core; the double pendulum reference at $h=10^{-4}$ over
$T=3$ took $1.6$~h, each of the other scripts between a minute and about an hour.
